% --- Archon physics formalization source begin ---
% archon:physics
% archon:covers PhyXMiniProblems/problem_phyx_mini_0841.lean
% archon:source-report reports/phyx_mini/problem_phyx_mini_0841.source.json

\chapter{Physics problem phyx\_mini\_0841}
\label{ch:PhyXMiniProblems_problem_phyx_mini_0841}

\paragraph{Problem source.}
\#\# Physical scenario

The cube contains negative charge. The electric field is constant over each face of the cube.

\#\# Question

What strength must this field exceed?

\#\# Answer choices

- A: A: -4N/C

- B: B: -15N/C

- C: C: -5.6N/C

- D: D: -5N/C

\#\# Figure caption (auxiliary; use the image as primary evidence)

The image depicts a three-dimensional cube with several arrows indicating field strengths around it. The cube is drawn with some faces shown as solid lines and others as dashed lines, giving a perspective view.

There are six arrows with labeled magnitudes, likely representing electric field vectors in newtons per coulomb (N/C):

- One arrow points downward from the top face of the cube with a magnitude of 15 N/C.
- One arrow points upward from the bottom face with a magnitude of 10 N/C.
- One arrow points to the left from the left face with a magnitude of 10 N/C.
- One arrow points to the right from the right face with a magnitude of 20 N/C.
- One arrow points outward toward the viewer from the front face with a magnitude of 20 N/C.
- One arrow points inward from the back face, not labeled but implied similar to the right arrow.

The text on the image reads "Field strengths in N/C," indicating the magnitudes are given in newtons per coulomb.

\#\# Dataset metadata

Electrostatics; Physical Model Grounding Reasoning, Multi-Formula Reasoning

\paragraph{Recorded answer/context.}
D: D: -5N/C

\paragraph{Figure/image path.}
/root/proposal\_for\_physic/science-mango/phyx\_mini\_run/phyx\_data/test\_image/841.png

\paragraph{Formalization target.}
create a compiling Lean file with sorry bodies at `PhyXMiniProblems/problem\_phyx\_mini\_0841.lean`. The Lean declarations must preserve the physical quantities, dimensions or dimensional roles, figure labels, governing-law hypotheses, and final relation expressed by this problem.
Use Mathlib/Physlib names found through LeanExplore where available. If a physics API is missing, introduce faithful local abstractions rather than scalar placeholder aliases.

\begin{theorem}[Physics formalization target]
\label{thm:physics:phyx_mini_0841:target}
\lean{PhyXMiniProblems.ProblemPhyXMini0841.unknownBackFaceFieldMustBeBelowNegativeFive}
\uses{def:physics:phyx-mini-0841:phyxminiproblems-problemphyxmini0841-fieldcomponentinnewtonspercoulomb, def:physics:phyx-mini-0841:phyxminiproblems-problemphyxmini0841-chargedcubesetup, def:physics:phyx-mini-0841:phyxminiproblems-problemphyxmini0841-matchesproblemandfigurereadouts, def:physics:phyx-mini-0841:phyxminiproblems-problemphyxmini0841-hasconstantfaceelectricfield, def:physics:phyx-mini-0841:phyxminiproblems-problemphyxmini0841-figurerepresentssurfacecomponents, def:physics:phyx-mini-0841:phyxminiproblems-problemphyxmini0841-hasphysicalcubeparameters, def:physics:phyx-mini-0841:phyxminiproblems-problemphyxmini0841-satisfiescubesurfacefluxlaws}
The assigned autoformalize agent should translate this physics problem into Lean declarations in the covered file, with theorem and lemma proof bodies written as `by sorry`.
\end{theorem}
\begin{proof}
This is an autoformalization task, not a proof task. Produce statements that can later be proved by the physics prover without weakening the physical model.
\end{proof}

% --- Archon declaration topology begin ---
\section{Declaration topology}
\begin{definition}[electric Field Strength Dimension]
  \label{def:physics:phyx-mini-0841:phyxminiproblems-problemphyxmini0841-electricfieldstrengthdimension}
  \lean{PhyXMiniProblems.ProblemPhyXMini0841.electricFieldStrengthDimension}
  The dimension M L T⁻² C⁻¹ of electric-field strength
\end{definition}
\begin{proof}
  This is the declared model definition of electric Field Strength Dimension.
\end{proof}
\begin{definition}[area Dimension]
  \label{def:physics:phyx-mini-0841:phyxminiproblems-problemphyxmini0841-areadimension}
  \lean{PhyXMiniProblems.ProblemPhyXMini0841.areaDimension}
  The dimension L² of a surface area
\end{definition}
\begin{proof}
  This is the declared model definition of area Dimension.
\end{proof}
\begin{definition}[electric Flux Dimension]
  \label{def:physics:phyx-mini-0841:phyxminiproblems-problemphyxmini0841-electricfluxdimension}
  \lean{PhyXMiniProblems.ProblemPhyXMini0841.electricFluxDimension}
  The dimension M L³ T⁻² C⁻¹ of electric flux
\end{definition}
\begin{proof}
  This is the declared model definition of electric Flux Dimension.
\end{proof}
\begin{definition}[electric Permittivity Dimension]
  \label{def:physics:phyx-mini-0841:phyxminiproblems-problemphyxmini0841-electricpermittivitydimension}
  \lean{PhyXMiniProblems.ProblemPhyXMini0841.electricPermittivityDimension}
  The dimension C² T² M⁻¹ L⁻³ of electric permittivity
\end{definition}
\begin{proof}
  This is the declared model definition of electric Permittivity Dimension.
\end{proof}
\begin{definition}[Length Quantity]
  \label{def:physics:phyx-mini-0841:phyxminiproblems-problemphyxmini0841-lengthquantity}
  \lean{PhyXMiniProblems.ProblemPhyXMini0841.LengthQuantity}
  A nonnegative unit-independent physical length
\end{definition}
\begin{proof}
  This is the declared model definition of Length Quantity.
\end{proof}
\begin{definition}[Area Quantity]
  \label{def:physics:phyx-mini-0841:phyxminiproblems-problemphyxmini0841-areaquantity}
  \lean{PhyXMiniProblems.ProblemPhyXMini0841.AreaQuantity}
  \uses{def:physics:phyx-mini-0841:phyxminiproblems-problemphyxmini0841-areadimension}
  A nonnegative unit-independent physical area
\end{definition}
\begin{proof}
  This is the declared model definition of Area Quantity.
\end{proof}
\begin{definition}[Signed Charge Quantity]
  \label{def:physics:phyx-mini-0841:phyxminiproblems-problemphyxmini0841-signedchargequantity}
  \lean{PhyXMiniProblems.ProblemPhyXMini0841.SignedChargeQuantity}
  A signed electric charge
\end{definition}
\begin{proof}
  This is the declared model definition of Signed Charge Quantity.
\end{proof}
\begin{definition}[Electric Field Magnitude Quantity]
  \label{def:physics:phyx-mini-0841:phyxminiproblems-problemphyxmini0841-electricfieldmagnitudequantity}
  \lean{PhyXMiniProblems.ProblemPhyXMini0841.ElectricFieldMagnitudeQuantity}
  \uses{def:physics:phyx-mini-0841:phyxminiproblems-problemphyxmini0841-electricfieldstrengthdimension}
  A nonnegative electric-field magnitude, as printed next to an arrow
\end{definition}
\begin{proof}
  This is the declared model definition of Electric Field Magnitude Quantity.
\end{proof}
\begin{definition}[Signed Electric Field Component Quantity]
  \label{def:physics:phyx-mini-0841:phyxminiproblems-problemphyxmini0841-signedelectricfieldcomponentquantity}
  \lean{PhyXMiniProblems.ProblemPhyXMini0841.SignedElectricFieldComponentQuantity}
  \uses{def:physics:phyx-mini-0841:phyxminiproblems-problemphyxmini0841-electricfieldstrengthdimension}
  An oriented outward-normal electric-field component
\end{definition}
\begin{proof}
  This is the declared model definition of Signed Electric Field Component Quantity.
\end{proof}
\begin{definition}[Electric Flux Quantity]
  \label{def:physics:phyx-mini-0841:phyxminiproblems-problemphyxmini0841-electricfluxquantity}
  \lean{PhyXMiniProblems.ProblemPhyXMini0841.ElectricFluxQuantity}
  \uses{def:physics:phyx-mini-0841:phyxminiproblems-problemphyxmini0841-electricfluxdimension}
  Signed electric flux through a face or a closed surface
\end{definition}
\begin{proof}
  This is the declared model definition of Electric Flux Quantity.
\end{proof}
\begin{definition}[Electric Permittivity Quantity]
  \label{def:physics:phyx-mini-0841:phyxminiproblems-problemphyxmini0841-electricpermittivityquantity}
  \lean{PhyXMiniProblems.ProblemPhyXMini0841.ElectricPermittivityQuantity}
  \uses{def:physics:phyx-mini-0841:phyxminiproblems-problemphyxmini0841-electricpermittivitydimension}
  A nonnegative electric permittivity
\end{definition}
\begin{proof}
  This is the declared model definition of Electric Permittivity Quantity.
\end{proof}
\begin{definition}[length In Meters]
  \label{def:physics:phyx-mini-0841:phyxminiproblems-problemphyxmini0841-lengthinmeters}
  \lean{PhyXMiniProblems.ProblemPhyXMini0841.lengthInMeters}
  \uses{def:physics:phyx-mini-0841:phyxminiproblems-problemphyxmini0841-lengthquantity}
  SI readout of a length, in metres
\end{definition}
\begin{proof}
  This is the declared model definition of length In Meters.
\end{proof}
\begin{definition}[area In Square Meters]
  \label{def:physics:phyx-mini-0841:phyxminiproblems-problemphyxmini0841-areainsquaremeters}
  \lean{PhyXMiniProblems.ProblemPhyXMini0841.areaInSquareMeters}
  \uses{def:physics:phyx-mini-0841:phyxminiproblems-problemphyxmini0841-areaquantity}
  SI readout of an area, in square metres
\end{definition}
\begin{proof}
  This is the declared model definition of area In Square Meters.
\end{proof}
\begin{definition}[charge In Coulombs]
  \label{def:physics:phyx-mini-0841:phyxminiproblems-problemphyxmini0841-chargeincoulombs}
  \lean{PhyXMiniProblems.ProblemPhyXMini0841.chargeInCoulombs}
  \uses{def:physics:phyx-mini-0841:phyxminiproblems-problemphyxmini0841-signedchargequantity}
  SI readout of signed charge, in coulombs
\end{definition}
\begin{proof}
  This is the declared model definition of charge In Coulombs.
\end{proof}
\begin{definition}[field Magnitude In Newtons Per Coulomb]
  \label{def:physics:phyx-mini-0841:phyxminiproblems-problemphyxmini0841-fieldmagnitudeinnewtonspercoulomb}
  \lean{PhyXMiniProblems.ProblemPhyXMini0841.fieldMagnitudeInNewtonsPerCoulomb}
  \uses{def:physics:phyx-mini-0841:phyxminiproblems-problemphyxmini0841-electricfieldmagnitudequantity}
  SI readout of a nonnegative field magnitude, in newtons per coulomb
\end{definition}
\begin{proof}
  This is the declared model definition of field Magnitude In Newtons Per Coulomb.
\end{proof}
\begin{definition}[field Component In Newtons Per Coulomb]
  \label{def:physics:phyx-mini-0841:phyxminiproblems-problemphyxmini0841-fieldcomponentinnewtonspercoulomb}
  \lean{PhyXMiniProblems.ProblemPhyXMini0841.fieldComponentInNewtonsPerCoulomb}
  \uses{def:physics:phyx-mini-0841:phyxminiproblems-problemphyxmini0841-signedelectricfieldcomponentquantity}
  SI readout of a signed normal field component, in newtons per coulomb
\end{definition}
\begin{proof}
  This is the declared model definition of field Component In Newtons Per Coulomb.
\end{proof}
\begin{definition}[electric Flux In Newton Square Meters Per Coulomb]
  \label{def:physics:phyx-mini-0841:phyxminiproblems-problemphyxmini0841-electricfluxinnewtonsquaremeterspercoulomb}
  \lean{PhyXMiniProblems.ProblemPhyXMini0841.electricFluxInNewtonSquareMetersPerCoulomb}
  \uses{def:physics:phyx-mini-0841:phyxminiproblems-problemphyxmini0841-electricfluxquantity}
  SI readout of electric flux, in newton-square-metres per coulomb
\end{definition}
\begin{proof}
  This is the declared model definition of electric Flux In Newton Square Meters Per Coulomb.
\end{proof}
\begin{definition}[permittivity In SI]
  \label{def:physics:phyx-mini-0841:phyxminiproblems-problemphyxmini0841-permittivityinsi}
  \lean{PhyXMiniProblems.ProblemPhyXMini0841.permittivityInSI}
  \uses{def:physics:phyx-mini-0841:phyxminiproblems-problemphyxmini0841-electricpermittivityquantity}
  SI readout of permittivity, in coulombs squared per newton-square-metre
\end{definition}
\begin{proof}
  This is the declared model definition of permittivity In SI.
\end{proof}
\begin{definition}[Cube Face]
  \label{def:physics:phyx-mini-0841:phyxminiproblems-problemphyxmini0841-cubeface}
  \lean{PhyXMiniProblems.ProblemPhyXMini0841.CubeFace}
  The six faces of the Gaussian cube
\end{definition}
\begin{proof}
  This is the declared model definition of Cube Face.
\end{proof}
\begin{definition}[Normal Orientation]
  \label{def:physics:phyx-mini-0841:phyxminiproblems-problemphyxmini0841-normalorientation}
  \lean{PhyXMiniProblems.ProblemPhyXMini0841.NormalOrientation}
  Orientation of a pictured arrow relative to a face's outward normal
\end{definition}
\begin{proof}
  This is the declared model definition of Normal Orientation.
\end{proof}
\begin{definition}[outward Sign]
  \label{def:physics:phyx-mini-0841:phyxminiproblems-problemphyxmini0841-normalorientation-outwardsign}
  \lean{PhyXMiniProblems.ProblemPhyXMini0841.NormalOrientation.outwardSign}
  \uses{def:physics:phyx-mini-0841:phyxminiproblems-problemphyxmini0841-normalorientation}
  The sign contributed by an arrow to the outward-normal component
\end{definition}
\begin{proof}
  This is the declared model definition of outward Sign.
\end{proof}
\begin{definition}[Answer Choice]
  \label{def:physics:phyx-mini-0841:phyxminiproblems-problemphyxmini0841-answerchoice}
  \lean{PhyXMiniProblems.ProblemPhyXMini0841.AnswerChoice}
  Answer labels recorded in the source problem
\end{definition}
\begin{proof}
  This is the declared model definition of Answer Choice.
\end{proof}
\begin{definition}[answer Choice Readout In Newtons Per Coulomb]
  \label{def:physics:phyx-mini-0841:phyxminiproblems-problemphyxmini0841-answerchoicereadoutinnewtonspercoulomb}
  \lean{PhyXMiniProblems.ProblemPhyXMini0841.answerChoiceReadoutInNewtonsPerCoulomb}
  \uses{def:physics:phyx-mini-0841:phyxminiproblems-problemphyxmini0841-answerchoice}
  The signed N/C threshold printed for each answer choice
\end{definition}
\begin{proof}
  This is the declared model definition of answer Choice Readout In Newtons Per Coulomb.
\end{proof}
\begin{definition}[Cube Electric Field Figure]
  \label{def:physics:phyx-mini-0841:phyxminiproblems-problemphyxmini0841-cubeelectricfieldfigure}
  \lean{PhyXMiniProblems.ProblemPhyXMini0841.CubeElectricFieldFigure}
  \uses{def:physics:phyx-mini-0841:phyxminiproblems-problemphyxmini0841-electricfieldmagnitudequantity, def:physics:phyx-mini-0841:phyxminiproblems-problemphyxmini0841-cubeface, def:physics:phyx-mini-0841:phyxminiproblems-problemphyxmini0841-normalorientation}
  Typed content of the supplied raster image. Magnitude labels are physical quantities; their numerical readings and arrow directions are imposed by the separate MatchesProblemAndFigureReadouts predicate
\end{definition}
\begin{proof}
  This is the declared model definition of Cube Electric Field Figure.
\end{proof}
\begin{definition}[Has Labeled Arrow]
  \label{def:physics:phyx-mini-0841:phyxminiproblems-problemphyxmini0841-cubeelectricfieldfigure-haslabeledarrow}
  \lean{PhyXMiniProblems.ProblemPhyXMini0841.CubeElectricFieldFigure.HasLabeledArrow}
  \uses{def:physics:phyx-mini-0841:phyxminiproblems-problemphyxmini0841-fieldmagnitudeinnewtonspercoulomb, def:physics:phyx-mini-0841:phyxminiproblems-problemphyxmini0841-cubeface, def:physics:phyx-mini-0841:phyxminiproblems-problemphyxmini0841-normalorientation, def:physics:phyx-mini-0841:phyxminiproblems-problemphyxmini0841-cubeelectricfieldfigure}
  One pictured face has an arrow of the stated orientation and SI magnitude
\end{definition}
\begin{proof}
  This is the declared model definition of Has Labeled Arrow.
\end{proof}
\begin{definition}[Charged Cube Setup]
  \label{def:physics:phyx-mini-0841:phyxminiproblems-problemphyxmini0841-chargedcubesetup}
  \lean{PhyXMiniProblems.ProblemPhyXMini0841.ChargedCubeSetup}
  \uses{def:physics:phyx-mini-0841:phyxminiproblems-problemphyxmini0841-lengthquantity, def:physics:phyx-mini-0841:phyxminiproblems-problemphyxmini0841-areaquantity, def:physics:phyx-mini-0841:phyxminiproblems-problemphyxmini0841-signedchargequantity, def:physics:phyx-mini-0841:phyxminiproblems-problemphyxmini0841-signedelectricfieldcomponentquantity, def:physics:phyx-mini-0841:phyxminiproblems-problemphyxmini0841-electricfluxquantity, def:physics:phyx-mini-0841:phyxminiproblems-problemphyxmini0841-electricpermittivityquantity, def:physics:phyx-mini-0841:phyxminiproblems-problemphyxmini0841-cubeface, def:physics:phyx-mini-0841:phyxminiproblems-problemphyxmini0841-cubeelectricfieldfigure}
  The cube, its ambient field, the derived normal components and fluxes, and the primary figure. No field component is defined from the desired -5 N/C threshold
\end{definition}
\begin{proof}
  This is the declared model definition of Charged Cube Setup.
\end{proof}
\begin{definition}[Matches Problem And Figure Readouts]
  \label{def:physics:phyx-mini-0841:phyxminiproblems-problemphyxmini0841-matchesproblemandfigurereadouts}
  \lean{PhyXMiniProblems.ProblemPhyXMini0841.MatchesProblemAndFigureReadouts}
  \uses{def:physics:phyx-mini-0841:phyxminiproblems-problemphyxmini0841-chargedcubesetup}
  Source-text and primary-image evidence, excluding the unknown back field
\end{definition}
\begin{proof}
  This is the declared model definition of Matches Problem And Figure Readouts.
\end{proof}
\begin{definition}[Has Constant Face Electric Field]
  \label{def:physics:phyx-mini-0841:phyxminiproblems-problemphyxmini0841-hasconstantfaceelectricfield}
  \lean{PhyXMiniProblems.ProblemPhyXMini0841.HasConstantFaceElectricField}
  \uses{def:physics:phyx-mini-0841:phyxminiproblems-problemphyxmini0841-fieldcomponentinnewtonspercoulomb, def:physics:phyx-mini-0841:phyxminiproblems-problemphyxmini0841-chargedcubesetup}
  The physical field is constant on each face, and the signed dimensionful component attached to a face is its projection onto the outward unit normal
\end{definition}
\begin{proof}
  This is the declared model definition of Has Constant Face Electric Field.
\end{proof}
\begin{definition}[Figure Represents Surface Components]
  \label{def:physics:phyx-mini-0841:phyxminiproblems-problemphyxmini0841-figurerepresentssurfacecomponents}
  \lean{PhyXMiniProblems.ProblemPhyXMini0841.FigureRepresentsSurfaceComponents}
  \uses{def:physics:phyx-mini-0841:phyxminiproblems-problemphyxmini0841-fieldmagnitudeinnewtonspercoulomb, def:physics:phyx-mini-0841:phyxminiproblems-problemphyxmini0841-fieldcomponentinnewtonspercoulomb, def:physics:phyx-mini-0841:phyxminiproblems-problemphyxmini0841-chargedcubesetup}
  The image's unsigned magnitude and arrow orientation determine a known face's signed outward-normal component. This applies uniformly to any labeled face and contains no condition on the unlabeled back face
\end{definition}
\begin{proof}
  This is the declared model definition of Figure Represents Surface Components.
\end{proof}
\begin{definition}[Has Physical Cube Parameters]
  \label{def:physics:phyx-mini-0841:phyxminiproblems-problemphyxmini0841-hasphysicalcubeparameters}
  \lean{PhyXMiniProblems.ProblemPhyXMini0841.HasPhysicalCubeParameters}
  \uses{def:physics:phyx-mini-0841:phyxminiproblems-problemphyxmini0841-lengthinmeters, def:physics:phyx-mini-0841:phyxminiproblems-problemphyxmini0841-areainsquaremeters, def:physics:phyx-mini-0841:phyxminiproblems-problemphyxmini0841-chargeincoulombs, def:physics:phyx-mini-0841:phyxminiproblems-problemphyxmini0841-permittivityinsi, def:physics:phyx-mini-0841:phyxminiproblems-problemphyxmini0841-chargedcubesetup}
  Positivity and charge-sign facts for the physical cube
\end{definition}
\begin{proof}
  This is the declared model definition of Has Physical Cube Parameters.
\end{proof}
\begin{definition}[Satisfies Cube Surface Flux Laws]
  \label{def:physics:phyx-mini-0841:phyxminiproblems-problemphyxmini0841-satisfiescubesurfacefluxlaws}
  \lean{PhyXMiniProblems.ProblemPhyXMini0841.SatisfiesCubeSurfaceFluxLaws}
  \uses{def:physics:phyx-mini-0841:phyxminiproblems-problemphyxmini0841-lengthinmeters, def:physics:phyx-mini-0841:phyxminiproblems-problemphyxmini0841-areainsquaremeters, def:physics:phyx-mini-0841:phyxminiproblems-problemphyxmini0841-chargeincoulombs, def:physics:phyx-mini-0841:phyxminiproblems-problemphyxmini0841-fieldcomponentinnewtonspercoulomb, def:physics:phyx-mini-0841:phyxminiproblems-problemphyxmini0841-electricfluxinnewtonsquaremeterspercoulomb, def:physics:phyx-mini-0841:phyxminiproblems-problemphyxmini0841-permittivityinsi, def:physics:phyx-mini-0841:phyxminiproblems-problemphyxmini0841-cubeface, def:physics:phyx-mini-0841:phyxminiproblems-problemphyxmini0841-chargedcubesetup}
  Equal-area cube geometry, constant-field face flux, flux additivity, and the integral form of Gauss's law. The last equation is the local abstraction needed because Physlib's available Gauss-law theorem is pointwise and differential rather than a closed-surface integral theorem
\end{definition}
\begin{proof}
  This is the declared model definition of Satisfies Cube Surface Flux Laws.
\end{proof}
\begin{lemma}[displayed Face Components Sum To Five]
  \label{lem:physics:phyx-mini-0841:phyxminiproblems-problemphyxmini0841-displayedfacecomponentssumtofive}
  \lean{PhyXMiniProblems.ProblemPhyXMini0841.displayedFaceComponentsSumToFive}
  \uses{def:physics:phyx-mini-0841:phyxminiproblems-problemphyxmini0841-fieldcomponentinnewtonspercoulomb, def:physics:phyx-mini-0841:phyxminiproblems-problemphyxmini0841-chargedcubesetup, def:physics:phyx-mini-0841:phyxminiproblems-problemphyxmini0841-matchesproblemandfigurereadouts, def:physics:phyx-mini-0841:phyxminiproblems-problemphyxmini0841-figurerepresentssurfacecomponents}
  The five displayed signed components have outward-normal sum 5 N/C
\end{lemma}
\begin{proof}
  The conclusion follows from the governing relations and figure readouts named in the statement of displayed Face Components Sum To Five.
\end{proof}
\begin{lemma}[total Outward Flux Is Negative]
  \label{lem:physics:phyx-mini-0841:phyxminiproblems-problemphyxmini0841-totaloutwardfluxisnegative}
  \lean{PhyXMiniProblems.ProblemPhyXMini0841.totalOutwardFluxIsNegative}
  \uses{def:physics:phyx-mini-0841:phyxminiproblems-problemphyxmini0841-electricfluxinnewtonsquaremeterspercoulomb, def:physics:phyx-mini-0841:phyxminiproblems-problemphyxmini0841-chargedcubesetup, def:physics:phyx-mini-0841:phyxminiproblems-problemphyxmini0841-hasphysicalcubeparameters, def:physics:phyx-mini-0841:phyxminiproblems-problemphyxmini0841-satisfiescubesurfacefluxlaws}
  Negative enclosed charge and positive permittivity make the total flux negative
\end{lemma}
\begin{proof}
  The conclusion follows from the governing relations and figure readouts named in the statement of total Outward Flux Is Negative.
\end{proof}
% --- Archon declaration topology end ---
% --- Archon physics formalization source end ---
